% ============================================================
\setpartpreamble{\partintro
  {Custom SQL expressions, aggregates, run-time metadata and scaffolding: the tools for very
  complex plugins.}
  {
    \item Chapter~\ref{ch:custom-sql-expression-nodes}: Custom SQL expression nodes
    \item Chapter~\ref{ch:aggregate-translators}: Aggregate translators
    \item Chapter~\ref{ch:property-level-fluent-api-and}: Property-level Fluent API and
    runtime annotations
    \item Chapter~\ref{ch:scaffolding-in-practice}: Scaffolding in practice
  }}
\part{Deep extensions}\label{part:deep-extensions}

% ------------------------------------------------------------
\chapter{Custom SQL expression nodes}\label{ch:custom-sql-expression-nodes}
\versionrisk{volatile}

Chapter~\ref{ch:translating-methods-and-members} builds SQL from existing building blocks:
functions, constants, casts. That is enough for most plugins. A \textbf{custom expression
node} (a subclass of \code{Sql\allowbreak{}Expression}) becomes necessary when at least one
of the following applies:

\begin{datatable}{H{0.80}H{1.20}}
  \textbf{Reason} & \textbf{Example} \\
  \tablehead
  The syntax is not a function call & Operators such as \code{@@}, \code{<->},
  \code{IS DISTINCT FROM}, keywords such as \code{AT TIME ZONE} \\
  \tablerule
  The nullability semantics differ & \code{IS DISTINCT FROM} is never NULL, even when an
  operand is NULL \\
  \tablerule
  Other parts should recognize the node & Optimizations that only apply to this expression
  \\
\end{datatable}

The sample implements the null-safe comparison
\code{EF.\allowbreak{}Functions.\allowbreak{}Is\allowbreak{}Not\allowbreak{}Distinct\allowbreak{}From(\allowbreak{}a,\allowbreak{}
b)} and \code{Is\allowbreak{}Distinct\allowbreak{}From(\allowbreak{}a,\allowbreak{} b)}
in a small package of its own, \code{Acme.\allowbreak{}Entity\allowbreak{}Framework\allowbreak{}Core.\allowbreak{}Functions.\allowbreak{}Postgre\allowbreak{}SQL}, activated with \code{Use\allowbreak{}Npgsql(\allowbreak{}cs,\allowbreak{} o => o.\allowbreak{}Use\allowbreak{}Acme\allowbreak{}Functions(\allowbreak{}))}.
Npgsql does not offer the comparison as a function; it meets all three reasons.

\begin{note}[A teaching example, deliberately not about tenancy]
  The comparison has nothing to do with multi-tenancy. It lives in a sample package of its own
  because a custom node forces a plugin to replace two exclusive provider services, the SQL
  generator and the nullability processor. Up to edition~1.1 the tenancy plugin did so for this
  example alone, which made it harder to combine with other plugins. A tenancy plugin needs no
  custom node; function composition (chapter~\ref{ch:translating-methods-and-members}) is
  enough. The tests run both plugins side by side, and the pgvector example at the end of this
  chapter adds a third one that derives from the functions plugin's generator.
\end{note}

\section{Anatomy of a node}

\begin{datatable}{H{0.78}H{1.26}H{0.96}}
  \textbf{Member} & \textbf{Purpose} & \textbf{Consequence if missing} \\
  \tablehead
  \code{Visit\allowbreak{}Children} / \code{Update} & Visitors can replace children;
  unchanged children yield the same instance & EF throws at the first visitor \\
  \tablerule
  \code{Equals} / \code{Get\allowbreak{}Hash\allowbreak{}Code} & Structural comparison for
  the query cache and optimizations & Cache misses, wrong deduplication \\
  \tablerule
  \code{Print} & Readable representation in DebugView and error messages & Unclear errors \\
  \tablerule
  \code{Quote} & C\# code for precompiled queries and NativeAOT & No AOT support \\
  \tablerule
  SQL generator & Writes the node as SQL & \textbf{Silently} missing SQL \\
  \tablerule
  Nullability processor & Knows the NULL semantics & Exception during translation \\
\end{datatable}

The last two rows also apply to combinations with other plugins: if a second plugin replaces
the generator or the processor, the node's SQL is silently missing or the query throws.
\code{Composition\allowbreak{}Matrix\allowbreak{}Tests} shows both
(chapter~\ref{ch:registering-services}).

% source: src/Acme.EntityFrameworkCore.Functions.PostgreSQL/Query/Expressions/DistinctFromExpression.cs
\begin{sourcecode}[Query/\allowbreak{}Expressions/\allowbreak{}Distinct\allowbreak{}From\allowbreak{}Expression.\allowbreak{}cs (\allowbreak{}shortened)]{csharp}
public class DistinctFromExpression : SqlExpression
{
    private static ConstructorInfo? _quotingConstructor;

    public DistinctFromExpression(SqlExpression left, SqlExpression right, bool negated,
        RelationalTypeMapping? typeMapping)
        : base(typeof(bool), typeMapping)
    {
        // ... ThrowIfNull ...
        Left = left;
        Right = right;
        IsNegated = negated;
    }

    public virtual SqlExpression Left { get; }
    public virtual SqlExpression Right { get; }
    public virtual bool IsNegated { get; }

    protected override Expression VisitChildren(ExpressionVisitor visitor)
        => Update((SqlExpression)visitor.Visit(Left), (SqlExpression)visitor.Visit(Right));

    public virtual DistinctFromExpression Update(SqlExpression left, SqlExpression right)
        => left == Left && right == Right
            ? this
            : new DistinctFromExpression(left, right, IsNegated, TypeMapping);

    public override Expression Quote()
        => New(
            _quotingConstructor ??= typeof(DistinctFromExpression).GetConstructor(
                [typeof(SqlExpression), typeof(SqlExpression), typeof(bool),
                 typeof(RelationalTypeMapping)])!,
            Left.Quote(), Right.Quote(), Constant(IsNegated),
            RelationalExpressionQuotingUtilities.QuoteTypeMapping(TypeMapping));

    protected override void Print(ExpressionPrinter expressionPrinter)
    {
        expressionPrinter.Visit(Left);
        expressionPrinter.Append(IsNegated ? " IS DISTINCT FROM " : " IS NOT DISTINCT FROM ");
        expressionPrinter.Visit(Right);
    }

    public override bool Equals(object? obj)
        => obj is DistinctFromExpression other
           && (ReferenceEquals(this, other)
               || base.Equals(other) && Left.Equals(other.Left)
                  && Right.Equals(other.Right) && IsNegated == other.IsNegated);

    public override int GetHashCode()
        => HashCode.Combine(base.GetHashCode(), Left, Right, IsNegated);
}
\end{sourcecode}

The translator creates the node and aligns the type mappings of both sides: if one side has
a mapping (a column, say), the other gets the same one. That way a constant becomes the
matching literal:

% source: src/Acme.EntityFrameworkCore.Functions.PostgreSQL/Query/DistinctFromTranslator.cs
\begin{sourcecode}[Query/\allowbreak{}Distinct\allowbreak{}From\allowbreak{}Translator.\allowbreak{}cs (\allowbreak{}excerpt)]{csharp}
var definition = method.GetGenericMethodDefinition();
if (definition != IsNotDistinctFromMethod && definition != IsDistinctFromMethod)
    return null;

var left = arguments[1];                        // arguments[0] is DbFunctions
var right = arguments[2];
var typeMapping = ExpressionExtensions.InferTypeMapping(left, right)
                  ?? _typeMappingSource.FindMapping(left.Type);

return new DistinctFromExpression(
    _sql.ApplyTypeMapping(left, typeMapping)!,
    _sql.ApplyTypeMapping(right, typeMapping)!,
    negated: definition == IsDistinctFromMethod,
    _typeMappingSource.FindMapping(typeof(bool)));
\end{sourcecode}

\section{Extending the SQL generator}

The \code{Query\allowbreak{}Sql\allowbreak{}Generator} writes SQL using a visitor. It hands
unknown nodes to
\code{Expression\allowbreak{}Visitor.\allowbreak{}Visit\allowbreak{}Extension}, which only
visits the children. Without an override there is no error, but \textbf{SQL without the
operator}. The functions plugin derives from the Npgsql generator:

% source: src/Acme.EntityFrameworkCore.Functions.PostgreSQL/Query/Internal/FunctionsNpgsqlQuerySqlGenerator.cs
\begin{sourcecode}[Query/\allowbreak{}Internal/\allowbreak{}Functions\allowbreak{}Npgsql\allowbreak{}Query\allowbreak{}Sql\allowbreak{}Generator.\allowbreak{}cs (\allowbreak{}shortened)]{csharp}
public class FunctionsNpgsqlQuerySqlGenerator : NpgsqlQuerySqlGenerator
{
    public FunctionsNpgsqlQuerySqlGenerator(QuerySqlGeneratorDependencies dependencies,
        IRelationalTypeMappingSource typeMappingSource,
        bool reverseNullOrderingEnabled, Version postgresVersion)
        : base(dependencies, typeMappingSource, reverseNullOrderingEnabled, postgresVersion)
    {
    }

    protected override Expression VisitExtension(Expression extensionExpression)
        => extensionExpression switch
        {
            DistinctFromExpression e => VisitDistinctFrom(e),
            _ => base.VisitExtension(extensionExpression),
        };

    protected virtual Expression VisitDistinctFrom(DistinctFromExpression expression)
    {
        VisitOperand(expression.Left);
        Sql.Append(expression.IsNegated ? " IS DISTINCT FROM " : " IS NOT DISTINCT FROM ");
        VisitOperand(expression.Right);
        return expression;
    }

    // The node as an operand of other operators: always in parentheses
    protected override bool RequiresParentheses(SqlExpression outerExpression,
        SqlExpression innerExpression)
        => innerExpression is DistinctFromExpression
           || base.RequiresParentheses(outerExpression, innerExpression);
}
\end{sourcecode}

Parentheses are needed in two directions. \code{Visit\allowbreak{}Operand} parenthesizes
compound operands \emph{inside} the node. For the node \emph{as} an operand, EF asks
\code{Requires\allowbreak{}Parentheses}, which only knows EF's own nodes. Without the override
the plugin wrote \code{a IS DISTINCT FROM 'x' <> a IS DISTINCT FROM 'y'}; PostgreSQL binds
\code{<>} tighter than \code{IS} and reports a syntax error.
\code{Deep\allowbreak{}Query\allowbreak{}Tests} checks the SQL,
\code{Advanced\allowbreak{}Tests} the query against PostgreSQL.

The generator is created through
\code{IQuery\allowbreak{}Sql\allowbreak{}Generator\allowbreak{}Factory}, a singleton single
service. The plugin replaces it with its own factory taking the same dependencies as
Npgsql's (\code{Query\allowbreak{}Sql\allowbreak{}Generator\allowbreak{}Dependencies}, type
mapping source, \code{INpgsql\allowbreak{}Singleton\allowbreak{}Options}).

\section{Extending the nullability processor}

After translation, EF optimizes NULL comparisons. To do so it must know for every node
whether it can become NULL. Unknown nodes end up in
\code{Visit\allowbreak{}Custom\allowbreak{}Sql\allowbreak{}Expression}, whose base
implementation throws an exception. The processor hangs off a chain of three classes, all of
which must be replaced:

% source: src/Acme.EntityFrameworkCore.Functions.PostgreSQL/Query/Internal/FunctionsNpgsqlSqlNullabilityProcessor.cs
\begin{sourcecode}[Query/\allowbreak{}Internal/\allowbreak{}Functions\allowbreak{}Npgsql\allowbreak{}Sql\allowbreak{}Nullability\allowbreak{}Processor.\allowbreak{}cs (\allowbreak{}shortened)]{csharp}
public class FunctionsNpgsqlSqlNullabilityProcessor : NpgsqlSqlNullabilityProcessor
{
    // Constructor: (…ProcessorDependencies, …ProcessorParameters) : base(...)

    protected override SqlExpression VisitCustomSqlExpression(
        SqlExpression sqlExpression, bool allowOptimizedExpansion, out bool nullable)
    {
        if (sqlExpression is DistinctFromExpression distinctFrom)
        {
            // Visit children (they can be optimized), their nullability does not matter
            var left = Visit(distinctFrom.Left, out _)!;
            var right = Visit(distinctFrom.Right, out _)!;

            nullable = false;                   // IS [NOT] DISTINCT FROM is never NULL
            return distinctFrom.Update(left, right);
        }

        return base.VisitCustomSqlExpression(
            sqlExpression, allowOptimizedExpansion, out nullable);
    }
}

public class FunctionsNpgsqlParameterBasedSqlProcessor : NpgsqlParameterBasedSqlProcessor
{
    protected override Expression ProcessSqlNullability(
        Expression selectExpression, ParametersCacheDecorator parametersDecorator)
        => new FunctionsNpgsqlSqlNullabilityProcessor(Dependencies, Parameters)
            .Process(selectExpression, parametersDecorator);
}

// + FunctionsNpgsqlParameterBasedSqlProcessorFactory
//   : IRelationalParameterBasedSqlProcessorFactory
\end{sourcecode}

% source: src/Acme.EntityFrameworkCore.Functions.PostgreSQL/Infrastructure/FunctionsOptionsExtension.cs
\begin{sourcecode}[Infrastructure/\allowbreak{}Functions\allowbreak{}Options\allowbreak{}Extension.\allowbreak{}cs › Apply\allowbreak{}Services (\allowbreak{}addition)]{csharp}
// 4) Deep extension: custom SQL expression node
services.Replace(ServiceDescriptor
    .Singleton<IQuerySqlGeneratorFactory, FunctionsNpgsqlQuerySqlGeneratorFactory>());
services.Replace(ServiceDescriptor
    .Scoped<IRelationalParameterBasedSqlProcessorFactory,
        FunctionsNpgsqlParameterBasedSqlProcessorFactory>());
\end{sourcecode}

\begin{important}[The price of custom nodes]
  Custom nodes require replacing two single services and deriving from internal Npgsql
  classes. Two plugins that both introduce custom nodes can only be combined if one derives
  from the other. Hence: function composition first; custom nodes only when one of the three
  reasons applies. And the replacement belongs prominently in the documentation.
\end{important}

\begin{note}[Experimental API]
  \code{Quote} and
  \code{Relational\allowbreak{}Expression\allowbreak{}Quoting\allowbreak{}Utilities} belong
  to the experimental API for precompiled queries. The compiler reports their use as EF9100;
  the sample suppresses it project-wide and documents this in
  \code{Directory.\allowbreak{}Build.\allowbreak{}props}.
\end{note}

% source: tests/Acme.EntityFrameworkCore.Tenancy.Tests/DeepQueryTests.cs
\begin{sourcecode}[tests/\allowbreak{}.\allowbreak{}.\allowbreak{}.\allowbreak{}/\allowbreak{}Deep\allowbreak{}Query\allowbreak{}Tests.\allowbreak{}cs (\allowbreak{}excerpt)]{csharp}
[Fact]
public void IsDistinctFrom_in_projection_passes_nullability_processor()
{
    // Without a custom nullability processor, EF would throw here
    using var ctx = Contexts.Create();

    var sql = ctx.Customers
        .Select(c => EF.Functions.IsDistinctFrom(c.Name, (string?)null))
        .ToQueryString();

    Assert.Contains("IS DISTINCT FROM", sql);
}
\end{sourcecode}

\section{Practical example: operators for vector search and pgvector}

Embeddings and similarity search bring operators into SQL that are not functions: pgvector
writes the L2 distance as \code{<->}, the cosine distance as
\code{<=>} and the negative inner product as
\code{<\#>}; full-text search uses \code{@@}.
\code{Sql\allowbreak{}Binary\allowbreak{}Expression} does not help here, because its operator
is an \code{Expression\allowbreak{}Type}, and there is none for
\code{<=>}. \code{Custom\allowbreak{}Operator\allowbreak{}Tests}
therefore shows
\code{EF.\allowbreak{}Functions.\allowbreak{}Cosine\allowbreak{}Distance(\allowbreak{}a,\allowbreak{}
b)} as a node of its own, \code{Vector\allowbreak{}Distance\allowbreak{}Expression}, with the
operator as a string, with translator, generator and nullability processor. It is test code:
an application uses \code{Pgvector.\allowbreak{}Entity\allowbreak{}Framework\allowbreak{}Core},
which brings the \code{vector} type, its type mapping and these translations. Npgsql itself has
\code{Pg\allowbreak{}Unknown\allowbreak{}Binary\allowbreak{}Expression} for such cases, a node
with an arbitrary operator and propagating nullability that needs no replaced generator or
processor. It lives in an \code{Internal} namespace, though (EF1001).

Three things decide whether such an operator produces correct SQL.

\textbf{Type mappings.} The translator gives both operands the same mapping, inferred from the
side that has one, usually the column. The parameter thus reaches the database with the
column's store type:

% source: tests/Acme.EntityFrameworkCore.Tenancy.Tests/CustomOperatorTests.cs
\begin{sourcecode}[tests/\allowbreak{}.\allowbreak{}.\allowbreak{}.\allowbreak{}/\allowbreak{}Custom\allowbreak{}Operator\allowbreak{}Tests.\allowbreak{}cs ›\allowbreak{} Translator]{csharp}
if (method != CosineDistanceMethod)
    return null;

var left = arguments[1];                        // arguments[0] is DbFunctions
var right = arguments[2];
var operandMapping = ExpressionExtensions.InferTypeMapping(left, right)
                     ?? _typeMappingSource.FindMapping(left.Type);

return new VectorDistanceExpression(
    _sql.ApplyTypeMapping(left, operandMapping)!,
    _sql.ApplyTypeMapping(right, operandMapping)!,
    "<=>",
    _typeMappingSource.FindMapping(typeof(double)));
\end{sourcecode}

In the test the column is a \code{float[]}, i.e. \code{real[]}. pgvector defines its operators
only for \code{vector}, though; against a real database this needs a type mapping with that
store type, as the package registers it. The result gets the mapping of \code{double}, i.e.
\code{double precision}.

\textbf{Precedence and parentheses.} EF does not know the precedence of a foreign operator.
PostgreSQL ranks user-defined operators between the arithmetic and the comparison operators:
\code{a <=> b * 2} means
\code{a <=> (\allowbreak{}b * 2)}. Without a countermeasure, the
generator wrote exactly that for
\code{Cosine\allowbreak{}Distance(\allowbreak{}.\allowbreak{}.\allowbreak{}.\allowbreak{}) * 2
+ 1}. The solution is the same two directions as above: parenthesize operands inside the node,
and always parenthesize the node as an operand through
\code{Requires\allowbreak{}Parentheses}. \code{Print} parenthesizes the same way, so that the
DebugView shows what the database reads.

% source: tests/Acme.EntityFrameworkCore.Tenancy.Tests/CustomOperatorTests.cs
\begin{sourcecode}[tests/\allowbreak{}.\allowbreak{}.\allowbreak{}.\allowbreak{}/\allowbreak{}Custom\allowbreak{}Operator\allowbreak{}Tests.\allowbreak{}cs ›\allowbreak{} Generator]{csharp}
public class VectorQuerySqlGenerator : FunctionsNpgsqlQuerySqlGenerator
{
    // ... constructor, VisitExtension, VisitOperand

    protected virtual Expression VisitVectorDistance(VectorDistanceExpression expression)
    {
        VisitOperand(expression.Left);
        Sql.Append($" {expression.Operator} ");
        VisitOperand(expression.Right);
        return expression;
    }

    // Without the override: "a <=> b * 2.0", read as a <=> (b * 2.0)
    protected override bool RequiresParentheses(SqlExpression outerExpression,
        SqlExpression innerExpression)
        => innerExpression is VectorDistanceExpression
           || base.RequiresParentheses(outerExpression, innerExpression);
}
\end{sourcecode}

The generator derives from the plugin's generator, not from Npgsql's: both replace the same
single service, and only this way do both nodes remain translatable. That is the price from the
box above, on a small scale.

\textbf{Nullability.} An ordinary operator is NULL as soon as an operand is NULL. The processor
reports that onwards:

% source: tests/Acme.EntityFrameworkCore.Tenancy.Tests/CustomOperatorTests.cs
\begin{sourcecode}[tests/\allowbreak{}.\allowbreak{}.\allowbreak{}.\allowbreak{}/\allowbreak{}Custom\allowbreak{}Operator\allowbreak{}Tests.\allowbreak{}cs ›\allowbreak{} Nullability]{csharp}
if (sqlExpression is VectorDistanceExpression distance)
{
    // NULL in, NULL out
    var left = Visit(distance.Left, out var leftNullable)!;
    var right = Visit(distance.Right, out var rightNullable)!;

    nullable = leftNullable || rightNullable;
    return distance.Update(left, right);
}
\end{sourcecode}

The report has visible consequences. For
\code{!(\allowbreak{}Cosine\allowbreak{}Distance(\allowbreak{}.\allowbreak{}.\allowbreak{}.\allowbreak{})
< 0.5)} EF writes \code{>= 0.5} for a required column, but a \code{CASE} for an optional one,
because \code{NOT (\allowbreak{}NULL < 0.5)} is NULL in SQL but \code{true} in C\#. Without the
processor, translation fails. The tests check the SQL of all cases:

% source: tests/Acme.EntityFrameworkCore.Tenancy.Tests/CustomOperatorTests.cs
\begin{sourcecode}[tests/\allowbreak{}.\allowbreak{}.\allowbreak{}.\allowbreak{}/\allowbreak{}Custom\allowbreak{}Operator\allowbreak{}Tests.\allowbreak{}cs (\allowbreak{}excerpt)]{csharp}
var sql = ctx.Documents
    .Select(d => EF.Functions.CosineDistance(d.Embedding, query) * 2 + 1)
    .ToQueryString();

Assert.Contains("SELECT (d.\"Embedding\" <=> @query) * 2.0 + 1.0", sql);
\end{sourcecode}

% ------------------------------------------------------------
\chapter{Aggregate translators}\label{ch:aggregate-translators}
\versionrisk{medium}

Aggregates such as \code{Sum} or \code{string\_\allowbreak{}agg} combine several rows,
typically after \code{Group\allowbreak{}By}. An
\code{IAggregate\allowbreak{}Method\allowbreak{}Call\allowbreak{}Translator} receives,
instead of a single argument, an \code{Enumerable\allowbreak{}Expression} that describes
\emph{what} is being aggregated:

\begin{datatable}{H{0.70}H{1.09}H{1.21}}
  \textbf{Property} & \textbf{Meaning} & \textbf{Origin in LINQ} \\
  \tablehead
  \code{Selector} & The aggregated expression & \code{g.\allowbreak{}Select(\allowbreak{}o
  => o.\allowbreak{}Total > 0)} \\
  \tablerule
  \code{Is\allowbreak{}Distinct} & Distinct values only &
  \code{g.\allowbreak{}Select(\allowbreak{}.\allowbreak{}.\allowbreak{}.\allowbreak{}).\allowbreak{}Distinct(\allowbreak{})}
  \\
  \tablerule
  \code{Predicate} & Filter before the aggregate & \code{g.\allowbreak{}Where(\allowbreak{}o
  => o.\allowbreak{}Id > 10)} \\
  \tablerule
  \code{Orderings} & Order (for order-dependent aggregates) &
  \code{g.\allowbreak{}Order\allowbreak{}By(\allowbreak{}.\allowbreak{}.\allowbreak{}.\allowbreak{})}
  \\
\end{datatable}

The functions sample package from chapter~\ref{ch:custom-sql-expression-nodes} adds
\code{EF.\allowbreak{}Functions.\allowbreak{}Every(\allowbreak{}values)} for PostgreSQL's
\code{bool\_\allowbreak{}and}; like the comparison, it is a teaching example unrelated to
tenancy:

% source: src/Acme.EntityFrameworkCore.Functions.PostgreSQL/Query/EveryAggregateTranslator.cs
\begin{sourcecode}[Query/\allowbreak{}Every\allowbreak{}Aggregate\allowbreak{}Translator.\allowbreak{}cs (\allowbreak{}excerpt)]{csharp}
public virtual SqlExpression? Translate(MethodInfo method, EnumerableExpression source,
    IReadOnlyList<SqlExpression> arguments,
    IDiagnosticsLogger<DbLoggerCategory.Query> logger)
{
    if (method != EveryMethod || source.Selector is not SqlExpression selector)
        return null;

    // Npgsql's factory implements DISTINCT, FILTER (WHERE ...) and ORDER BY in aggregates
    if (_sql is not NpgsqlSqlExpressionFactory npgsql)
        return null;

    var boolean = _typeMappingSource.FindMapping(typeof(bool))!;

    // bool_and returns NULL for an empty set, Enumerable.All returns true
    return _sql.Coalesce(
        npgsql.AggregateFunction("bool_and", [selector], source,
            nullable: true, argumentsPropagateNullability: [false],
            typeof(bool), boolean),
        _sql.Constant(true, boolean),
        boolean);
}
\end{sourcecode}

\begin{sourcecode}[Generated SQL for g.\allowbreak{}Where(\allowbreak{}o => o.\allowbreak{}Id > 10).\allowbreak{}Select(\allowbreak{}o => o.\allowbreak{}Total > 0)]{sql}
SELECT COALESCE(bool_and(o."Total" > 0.0) FILTER (WHERE o."Id" > 10), TRUE)
FROM "Orders" AS o
GROUP BY o."TenantId"
\end{sourcecode}

\begin{important}[Semantics before syntax]
  An aggregate translator must match the semantics of the .NET side, not just the database
  syntax. Typical pitfalls: empty sets (NULL instead of 0, true or an empty string), NULL
  values in the input, and predicates. Npgsql itself aligns \code{string.\allowbreak{}Join}
  with \code{string\_\allowbreak{}agg} using \code{COALESCE}; the sample follows this
  pattern.
\end{important}

Databases without a \code{FILTER} clause (such as SQL Server) emulate predicates with
\code{CASE WHEN predicate THEN value END}. A provider-neutral aggregate translator must know
both; that is why this one lives in the PostgreSQL package.

% ------------------------------------------------------------
\chapter{Property-level Fluent API and runtime annotations}\label{ch:property-level-fluent-api-and}
\versionrisk{volatile}

\section{Property builder}

So far the plugin assumes that the tenant property is called \code{Tenant\allowbreak{}Id}
and that \code{ITenant\allowbreak{}Owned} is implemented. With a Fluent API method on
\code{Property\allowbreak{}Builder} it works without either, for example with existing
models using a property \code{Mandant\allowbreak{}Id} or an explicit interface
implementation:

\begin{sourcecode}[Usage]{csharp}
modelBuilder.Entity<Contract>().Property(c => c.Mandant).IsTenantProperty();
\end{sourcecode}

% source: src/Acme.EntityFrameworkCore.Tenancy/Metadata/TenancyPropertyBuilderExtensions.cs
\begin{sourcecode}[Metadata/\allowbreak{}Tenancy\allowbreak{}Property\allowbreak{}Builder\allowbreak{}Extensions.\allowbreak{}cs (\allowbreak{}core package)]{csharp}
public static PropertyBuilder IsTenantProperty(
    this PropertyBuilder builder, bool isTenant = true)
{
    ArgumentNullException.ThrowIfNull(builder);

    if (isTenant)
        builder.Metadata.SetAnnotation(TenancyAnnotationNames.IsTenantProperty, true);
    else
        builder.Metadata.RemoveAnnotation(TenancyAnnotationNames.IsTenantProperty);

    return builder;
}

public static PropertyBuilder<TProperty> IsTenantProperty<TProperty>(
    this PropertyBuilder<TProperty> builder, bool isTenant = true)
    => (PropertyBuilder<TProperty>)IsTenantProperty((PropertyBuilder)builder, isTenant);

public static IConventionPropertyBuilder? IsTenantProperty(
    this IConventionPropertyBuilder builder,
    bool isTenant = true,
    bool fromDataAnnotation = false)
    => builder.HasAnnotation(TenancyAnnotationNames.IsTenantProperty,
           isTenant ? true : null, fromDataAnnotation) is null
        ? null
        : builder;
\end{sourcecode}

The method lives in the core package because it only sets metadata. The conventions of both
provider packages find the tenant property through the same core function:

\begin{sourcecode}[Metadata/\allowbreak{}Tenancy\allowbreak{}Entity\allowbreak{}Type\allowbreak{}Extensions.\allowbreak{}cs (\allowbreak{}core package,\allowbreak{} shortened, simplified)]{csharp}
public static IReadOnlyProperty? FindTenantPropertyCandidate(
    this IReadOnlyEntityType entityType)
{
    var marked = entityType.GetProperties()
        .Where(p => p[TenancyAnnotationNames.IsTenantProperty] is true)
        .ToList();

    if (marked.Count > 1)
        throw new InvalidOperationException("... several tenant properties ...");
    if (marked.Count == 1)
        return marked[0];

    return typeof(ITenantOwned).IsAssignableFrom(entityType.ClrType)
        ? entityType.FindProperty(nameof(ITenantOwned.TenantId))
        : null;
}
\end{sourcecode}

\begin{datatable}{H{0.92}H{1.10}H{0.98}}
  \textbf{Builder} & \textbf{Typical plugin methods} & \textbf{Note} \\
  \tablehead
  \code{Property\allowbreak{}Builder} / \code{<\allowbreak{}T>} & Column options, roles of
  properties & Here: \code{Is\allowbreak{}Tenant\allowbreak{}Property} \\
  \tablerule
  \code{Index\allowbreak{}Builder} / \code{<\allowbreak{}T>} & Index methods, operator
  classes, filters & Annotation on \code{Metadata} (IMutableIndex) \\
  \tablerule
  \code{Key\allowbreak{}Builder} & Key options (clustering) & As above \\
  \tablerule
  \code{Complex\allowbreak{}Type\allowbreak{}Property\allowbreak{}Builder} & Properties in
  complex types & Only offer where the option makes sense there
  (chapter~\ref{ch:complex-types-and-json}) \\
  \tablerule
  \code{Model\allowbreak{}Builder} & Model-wide options & Annotation on the model \\
\end{datatable}

\section{Runtime annotations}

The SaveChanges interceptor needs to know at run time which entities are isolated and what
their tenant property is called. The convention writes this into an annotation on the entity
type; the interceptor reads it from the \emph{runtime model}:

\begin{sourcecode}[Convention and core package (\allowbreak{}excerpt)]{csharp}
// Convention (design-time model)
entityType.Builder.HasAnnotation(TenancyAnnotationNames.TenantProperty, property.Name);

// Interceptor (runtime model)
public static string? GetTenantPropertyName(this IReadOnlyEntityType entityType)
    => entityType.GetRootType()[TenancyAnnotationNames.TenantProperty] as string;
\end{sourcecode}

Two questions decide whether this works. Does the annotation reach the runtime model? Yes:
\code{Runtime\allowbreak{}Model\allowbreak{}Convention} only removes EF's own design-time
annotations and keeps all others. Does it reach the compiled model? Yes: EF writes simple
values automatically (chapter~\ref{ch:compiled-models-trimming-and-nativeaot}). The custom
\code{ICSharp\allowbreak{}Runtime\allowbreak{}Annotation\allowbreak{}Code\allowbreak{}Generator}
therefore has the opposite job: it removes the RLS annotations that only migrations need:

\begin{sourcecode}[Design/\allowbreak{}Internal/\allowbreak{}Tenancy\allowbreak{}CSharp\allowbreak{}Runtime\allowbreak{}Annotation\allowbreak{}Code\allowbreak{}Generator.\allowbreak{}cs (\allowbreak{}shortened, simplified)]{csharp}
public class TenancyCSharpRuntimeAnnotationCodeGenerator
    : NpgsqlCSharpRuntimeAnnotationCodeGenerator
{
    // Constructor: (CSharpRuntimeAnnotationCodeGeneratorDependencies,
    //               RelationalCSharpRuntimeAnnotationCodeGeneratorDependencies) : base(...)

    public override void Generate(IEntityType entityType,
        CSharpRuntimeAnnotationCodeGeneratorParameters parameters)
    {
        if (!parameters.IsRuntime)
        {
            foreach (var key in parameters.Annotations.Keys
                         .Where(k => k.StartsWith(RlsAnnotationNames.Prefix,
                                                  StringComparison.Ordinal))
                         .ToList())
            {
                parameters.Annotations.Remove(key);
            }
        }

        base.Generate(entityType, parameters);
    }
}
\end{sourcecode}

When the compiled model is generated for NativeAOT
(\code{For\allowbreak{}Native\allowbreak{}Aot}), EF also writes the relational model. The
annotation provider has copied the RLS annotations onto the tables, so they appear a second
time there. The generator therefore applies the same filter in an override of
\code{Generate(\allowbreak{}ITable,\allowbreak{}
.\allowbreak{}.\allowbreak{}.\allowbreak{})}.

\begin{datatable}{H{1.36}H{0.88}H{0.88}H{0.88}}
  \textbf{Annotation} & \textbf{Design time} & \textbf{Run time} & \textbf{Compiled model} \\
  \tablehead
  \code{Acme:Rls:*} (policies) & yes (migrations) & yes, unused & filtered \\
  \tablerule
  \code{Acme:Tenancy:Tenant\allowbreak{}Property} & yes & yes (interceptor) & yes \\
  \tablerule
  \code{Acme:Tenancy:Is\allowbreak{}Tenant\allowbreak{}Property} & yes (convention) & yes,
  unused & yes (small) \\
\end{datatable}

\begin{note}[Two kinds of runtime annotations]
  In this book, “runtime annotation” means an ordinary annotation that is read at run time.
  EF also has a runtime annotation API in a narrower sense
  (\code{Add\allowbreak{}Runtime\allowbreak{}Annotation},
  \code{Get\allowbreak{}Or\allowbreak{}Add\allowbreak{}Runtime\allowbreak{}Annotation\allowbreak{}Value}
  on \code{IAnnotatable}). It stores values on the finished model only, for example derived
  data that should not be recomputed on every access. Such values never appear in the
  snapshot or in migrations. For configuration that a plugin sets while the model is built,
  the ordinary annotation is the right choice.
\end{note}

\begin{tip}
  For every annotation it should be clear who reads it and when. Annotations read only at
  design time do not belong in the compiled model; those read at run time must be in it and
  need a test against the generated model. The sample's compiled model test checks both.
\end{tip}

% ------------------------------------------------------------
\chapter{Scaffolding in practice}\label{ch:scaffolding-in-practice}
\versionrisk{volatile}

When scaffolding (\code{dotnet ef dbcontext scaffold}), EF reads the schema and generates
entities and a context. A plugin managing its own schema objects should recognize them in
the process. The route has three stations, two of which the plugin already knows:

\begin{flowdiagram}
  \flowstep{IDatabaseModelFactory}{Reads the schema into a DatabaseModel; the plugin
  attaches RLS data as annotations to the tables}
  \flowstep{Scaffolding model factory}{EF transfers table annotations to the entity types}
  \flowstep{IAnnotationCodeGenerator}{Writes them as HasRowLevelSecurity/HasPolicy
  (chapter~\ref{ch:design-time-services})}
\end{flowdiagram}

% source: src/Acme.EntityFrameworkCore.Tenancy.PostgreSQL/Scaffolding/TenancyDatabaseModelFactory.cs
\begin{sourcecode}[Scaffolding/\allowbreak{}Tenancy\allowbreak{}Database\allowbreak{}Model\allowbreak{}Factory.\allowbreak{}cs (\allowbreak{}shortened)]{csharp}
public class TenancyDatabaseModelFactory : NpgsqlDatabaseModelFactory
{
    public TenancyDatabaseModelFactory(
        IDiagnosticsLogger<DbLoggerCategory.Scaffolding> logger)
        : base(logger)
    {
    }

    public override DatabaseModel Create(
        DbConnection dbConnection, DatabaseModelFactoryOptions options)
    {
        var model = base.Create(dbConnection, options);
        // ... open the connection if needed ...
        ReadRowLevelSecurity(dbConnection, model);    // pg_class.relrowsecurity / ...force...
        ReadPolicies(dbConnection, model);            // pg_policies
        return model;
    }

    protected virtual void ReadPolicies(DbConnection connection, DatabaseModel model)
    {
        // SELECT schemaname, tablename, policyname, permissive, roles::text[], cmd,
        //        qual, with_check FROM pg_catalog.pg_policies
        // ...
        var policy = new RlsPolicy(
            Name: reader.GetString(2),
            Command: Enum.Parse<PolicyCommand>(reader.GetString(5), ignoreCase: true),
            Permissive: reader.GetString(3) == "PERMISSIVE",
            Roles: roles is ["public"] ? null : roles,    // "public" = default without TO
            Using: reader.IsDBNull(6) ? null : reader.GetString(6),
            WithCheck: reader.IsDBNull(7) ? null : reader.GetString(7));

        table[RlsAnnotationNames.PolicyPrefix + policy.Name] = policy.Serialize();
    }
}
\end{sourcecode}

The factory is registered like the code generator in the design-time services with
\code{Remove\allowbreak{}All}, but with \code{Add\allowbreak{}Scoped}: EF and Npgsql
register the database model factory as scoped, and it depends on the scaffolding logger,
which is scoped as well. The data format is the same as with the Fluent API; snapshot, migrations and code generator therefore work without
further changes.

\begin{note}[Normalized expressions]
  PostgreSQL stores policy expressions in normalized form, for example
  \code{(\allowbreak{}(\allowbreak{}\textquotedbl{}Tenant\allowbreak{}Id\textquotedbl{} =
  (\allowbreak{}NULLIF(\allowbreak{}current\_\allowbreak{}setting(\allowbreak{}'app.\allowbreak{}tenant\_\allowbreak{}id'::text,\allowbreak{}
  true),\allowbreak{} ''::text))::uuid))}. Scaffolded policies contain exactly this text.
  That is correct, but looks different from the form generated by the convention. Teams
  switching to migrations after scaffolding should unify the policies once, otherwise the
  first migration produces a DROP/CREATE with identical content.
\end{note}

% source: tests/Acme.EntityFrameworkCore.Tenancy.IntegrationTests/AdvancedTests.cs
\begin{sourcecode}[tests/\allowbreak{}.\allowbreak{}.\allowbreak{}.\allowbreak{}Integration\allowbreak{}Tests/\allowbreak{}Advanced\allowbreak{}Tests.\allowbreak{}cs (\allowbreak{}excerpt)]{csharp}
[Fact]
public void Scaffolding_reads_rls_and_policies()
{
    using var ctx = ContextFor(null);
    var factory = new TenancyDatabaseModelFactory(
        ctx.GetService<IDiagnosticsLogger<DbLoggerCategory.Scaffolding>>());

    var model = factory.Create(
        _fixture.OwnerConnectionString, new DatabaseModelFactoryOptions());
    var orders = model.Tables.Single(t => t.Name == "Orders");

    Assert.Equal(true, orders[RlsAnnotationNames.Enabled]);
    Assert.Equal(true, orders[RlsAnnotationNames.Forced]);
}
\end{sourcecode}
